\documentclass[11pt]{article}
\usepackage[margin=1in]{geometry}
\usepackage{times}
\usepackage{amsmath,amssymb}
\usepackage{graphicx}
\usepackage{booktabs}
\usepackage{array}
\usepackage{hyperref}
\usepackage{xcolor}
\usepackage{enumitem}
\usepackage{listings}
\usepackage{fancyhdr}
\usepackage{titlesec}

\hypersetup{
  colorlinks=true,
  linkcolor=blue,
  urlcolor=blue,
  citecolor=blue,
  pdftitle={moe-engine: A Fault-Tolerant Runtime for Hyperscale Mixture-of-Experts Training},
  pdfauthor={Min Htet Myet}
}

\lstset{
  basicstyle=\ttfamily\small,
  breaklines=true,
  columns=fullflexible,
  frame=single,
  backgroundcolor=\color{gray!5}
}

\pagestyle{fancy}
\setlength{\headheight}{14pt}
\fancyhf{}
\fancyhead[L]{moe-engine Preprint v3}
\fancyhead[R]{July 2026}
\fancyfoot[C]{\thepage}

\titlespacing*{\section}{0pt}{2.2ex plus 1ex minus .2ex}{1.5ex plus .2ex}
\titlespacing*{\subsection}{0pt}{1.8ex plus 1ex minus .2ex}{1ex plus .2ex}

\title{\textbf{moe-engine: A Fault-Tolerant Runtime\\for Hyperscale Mixture-of-Experts Training}\\
\large Version 3 Preprint}
\author{Min Htet Myet \\ \textit{Independent Researcher}}
\date{July 2026}

\begin{document}
\maketitle

\begin{center}
\url{https://github.com/Mattral/Composed-Mixture-of-Experts-Engine} \\
Zenodo (v1): \url{https://doi.org/10.5281/zenodo.20647577} \\
Zenodo (v2): \url{https://doi.org/10.5281/zenodo.20688837} \\
v3 in preparation
\end{center}

\begin{abstract}
\noindent
We present moe-engine, a research-grade infrastructure runtime for training
large Mixture-of-Experts (MoE) models under the realistic constraint of
continuous node failures at hyperscale (10K+ GPUs). The system combines a
high-performance fused Triton router, composable 4D parallelism (Data,
Expert, Tensor, Pipeline), strict mathematical invariants for correctness,
and elastic fault tolerance with automatic expert resharding.

This v3 preprint reports the two most significant developments since v2: (i)
the \emph{first sustained real-GPU validation} of the runtime on NVIDIA T4
hardware, which surfaced and fixed a Triton compilation defect invisible on
CPU-only continuous integration and produced measured throughput data
(80.1$\times$ GPU speedup over the CPU reference path at $N{=}4096$,
$H{=}2048$; 100\% pass rate on the storage-stall chaos scenario); and (ii) a
structural refactoring that replaced a 1{,}165-line monolithic module with
seven single-responsibility components, migrated configuration handling to a
statically validated schema, and added a Switch-Transformer-style expert
capacity enforcement mechanism. A related engineering finding is reported in
detail because of its methodological interest: a missing dependency
declaration caused the runtime's configuration validator to silently
degrade to a no-op across eight test cases in continuous integration for an
entire release cycle, a failure mode that is generalizable to any system
that ships an optional-dependency fallback for what is, in fact, safety-critical
logic. The test suite has grown from 148 to 348 passing cases, gaining
property-based testing and a mocked collective-communication backend that
exercises expert-parallel dispatch logic without a multi-process
environment.

We maintain the deliberately honest posture established in v1 and continued
in v2: sustained multi-node GPU benchmarks, a full resolution of the ~85\%
chaos-recovery pass rate, and hardware-scale validation of the newly added
expert-capacity mechanism at $E{=}256$ remain open items for v0.4, and are
reported as such rather than elided.
\end{abstract}

\section{Introduction}

Training large MoE models at true hyperscale exposes systems challenges
that go beyond raw throughput: continuous node failures, silent collective
and sharding bugs, misleading utilization metrics for sparse models, and
the need for seamless recovery with expert re-ownership. moe-engine was
designed from the ground up to address these challenges, and this is its
third preprint.

Version 1 described the initial architecture: the Triton router, 4D mesh
foundations, async two-tier checkpointing, and elastic recovery mechanisms.
Version 2 reported the completion of multi-process Pipeline Parallelism
with real inter-stage \texttt{send}/\texttt{recv} communication, a full
1F1B schedule, and fused Sequence Parallelism backward communication ---
but explicitly declined to present GPU throughput numbers, stating that
``sustained, reproducible cluster data is not yet available'' and that the
authors ``prefer to report real, reproducible results rather than
optimistic small-scale or illustrative numbers.''

This v3 preprint reports the resolution of that specific gap at single-GPU
scale: the system has now been run, measured, and validated on real CUDA
hardware for the first time. It also reports a substantial internal
refactoring undertaken in parallel, motivated by an external technical
review that identified the codebase's single-file distributed-systems
module as a maintainability risk disproportionate to the project's
research-grade maturity level. We report both the architecture-level
changes and a specific software-engineering failure mode discovered during
this work, because we believe the latter is of independent interest to
anyone building validated-configuration systems with optional-dependency
fallback paths.

\section{Background and Motivation}

Sections on MoE fundamentals, existing systems (Megatron-LM, DeepSpeed,
Switch Transformers, Tutel, X-MoE, MegaScale-MoE), and the hyperscale
failure model remain unchanged from v1/v2 and are omitted here for
brevity. The core motivation --- surviving continuous node death while
maintaining correctness --- is unchanged.

\section{System Architecture Overview}

The high-level architecture (training loop $\rightarrow$
\texttt{DistributedMoELayer} $\rightarrow$ \texttt{ElasticTrainerHarness}
$\rightarrow$ Telemetry) remains consistent with v1/v2. The major internal
change since v2 is structural rather than behavioral: the distributed
package, previously a single 1{,}165-line file
(\texttt{pkg/distributed/parallel\_mesh.py}) containing every parallelism
primitive, is now split into eight modules, each with one responsibility
(Table~\ref{tab:modules}). A 68-line backward-compatibility shim re-exports
every public symbol from the original location, so existing import
statements in downstream code and in the test suite continue to function
unmodified.

\begin{table}[h]
\centering
\small
\caption{Distributed package structure after the v0.3.2 refactor}
\label{tab:modules}
\begin{tabular}{@{}ll@{}}
\toprule
\textbf{Module} & \textbf{Responsibility} \\
\midrule
\texttt{mesh.py} & \texttt{ParallelTopology}, device mesh, comm streams \\
\texttt{tensor\_parallel.py} & Column/RowParallelLinear \\
\texttt{sequence\_parallel.py} & SP scatter/gather, fused all-gather \\
\texttt{expert\_parallel.py} & all-to-all dispatch/combine \\
\texttt{pipeline\_parallel.py} & \texttt{PipelineStage}, 1F1B schedule \\
\texttt{data\_parallel.py} & FSDP2 application \\
\texttt{router.py} & High-level router interface \\
\texttt{moe\_layer.py} & \texttt{DistributedMoELayer} orchestration \\
\texttt{parallel\_mesh.py} & Backward-compat.\ shim (68 lines) \\
\bottomrule
\end{tabular}
\end{table}

Two further architectural additions were made in this cycle. First, model
definitions were extracted from the training entrypoint into
\texttt{pkg/models/}, and a decorator-based registry
(\texttt{@register\_model}, \texttt{build\_model\_from\_config}) was added
so that the training loop dispatches on a configuration string rather than
importing a specific model class --- the pattern used by Megatron-LM and
torchTitan to decouple infrastructure code from model research code.
Second, the flat-dictionary configuration format used since v1 was replaced
with a statically validated schema (Section~\ref{sec:config}).

\section{Triton MoE Router Kernel}
\label{sec:triton}

The fused Triton router (forward + analytic backward) described in v1/v2
is architecturally unchanged: it continues to deliver single-HBM-pass
routing with strong numerical guarantees and forward-pass invariants. What
changed is that it has, for the first time, actually been exercised on GPU
hardware as part of this project's own validation process, and this
exposed a defect.

\subsection{The \texttt{constexpr} defect}

Both the forward and backward Triton kernels use \texttt{tl.static\_range}
to unroll the top-$K$ selection loop at compile time. This requires the
loop bound $K$ to be a Triton \texttt{constexpr} value. In the kernel
signature prior to this fix, $K$ was declared as an ordinary integer
parameter. On the CPU-only reference path used throughout v1 and v2's
continuous integration, this made no difference, because the reference
implementation is plain PyTorch and never invokes the Triton compiler. The
first time the kernel was compiled on real GPU hardware --- during the
work reported in this preprint --- it failed with a Triton
\texttt{CompilationError}, because \texttt{tl.static\_range} cannot unroll
a loop whose bound is not known at trace time.

We consider this defect noteworthy beyond its immediate fix (declaring $K$
as \texttt{constexpr} in both kernel signatures) because it is an instance
of a general failure mode in projects that maintain a CPU reference
fallback for hardware-specific code: the fallback path's very existence
allows a real hardware-only defect to survive an arbitrary number of CI
runs undetected, because CI never exercises the code path where the defect
lives. Two regression tests were added: one that statically asserts both
kernel signatures declare $K$ as \texttt{constexpr} (runs on CPU, catches
this specific class of regression without hardware), and one that
compiles and runs the kernel on GPU when available.

\subsection{Real-hardware measurements}

Section~\ref{sec:eval} reports the first real GPU throughput numbers for
this kernel, obtained on an NVIDIA T4 after the fix above.

\section{4D Parallelism and DistributedMoELayer}

\subsection{Device Mesh, Topology, and Sequence Parallelism}

\texttt{ParallelTopology} and device mesh construction are functionally
unchanged from v2; Sequence Parallelism was extracted from
\texttt{tensor\_parallel.py} into its own module
(\texttt{sequence\_parallel.py}) as part of the structural refactor, with
backward-compatible re-exports. All four mesh axes (DP, EP, TP, PP) remain
first-class.

\subsection{Expert Parallelism and Tensor Parallelism}

EP dispatch/combine with dedicated-stream overlap and the SwiGLU sharding
rule (both gate and up projections as \texttt{ColumnParallelLinear}) are
unchanged from v1/v2.

\subsection{Pipeline Parallelism}

The multi-process \texttt{PipelineStage} implementation reported in v2 ---
real \texttt{torch.distributed.send}/\texttt{recv} on a dedicated pipeline
process group, activation/gradient tagging by micro-batch index, and the
three-phase 1F1B schedule --- is unchanged in this release. It remains
validated by 2-rank \texttt{mp.spawn} tests and has not yet been exercised
under chaos conditions (Section~\ref{sec:limitations}).

\subsection{Sequence Parallelism Fusion}

The v2 fusion of Sequence Parallelism's backward communication with the
subsequent projection matmul (replacing an all-gather + matmul pattern with
a single all-reduce) is unchanged and remains validated by 2-rank
\texttt{mp.spawn} tests.

\subsection{Expert Capacity Enforcement (new in v3)}
\label{sec:capacity}

The single largest functional addition in this release is an optional,
opt-in expert-capacity enforcement mechanism, following the
first-come-first-served dropping policy used by Switch Transformer
\cite{fedus2021switch} and GShard \cite{lepikhin2020gshard}. Each expert
accepts at most $\lceil c \cdot N K / E \rceil$ tokens, where $c$ is a
configurable capacity factor, $N$ is the number of tokens, $K$ is the
number of active experts per token, and $E$ is the total expert count;
tokens beyond this budget receive zero combine weight for that slot rather
than being processed.

The implementation centers on a single primitive,
\texttt{\_cumcount(groups)}, which computes for each element its
zero-indexed position of appearance within its group, using a stable sort
followed by a \texttt{cummax} operation. This avoids both an explicit
Python loop over experts and a dependency on \texttt{scatter\_reduce}
variants whose availability differs across PyTorch versions. Capacity
enforcement is disabled by default (\texttt{capacity\_dropping=False}),
so no existing behavior changes unless a caller explicitly opts in; 22
dedicated tests cover the primitive, the mask construction, and full-layer
integration including gradient flow through non-dropped slots.

A new example configuration, \texttt{configs/large\_scale.yaml}, exercises
this mechanism at a fine-grained MoE scale ($E{=}256$, $K{=}8$, four times
the expert count of the existing production configuration) together with
the auxiliary load-balancing loss described below. This configuration has
been validated for correctness (forward/backward pass, capacity budget
arithmetic) at toy hidden dimensions; its throughput and memory behavior
at full scale on real GPU hardware has not yet been measured, and we do
not claim otherwise (Section~\ref{sec:limitations}).

\section{Configuration System}
\label{sec:config}

\subsection{Migration to a validated schema}

The flat-dictionary configuration format used since v1 provided no
validation: a value such as $\mathrm{top\_k} > \mathrm{num\_experts}$ would
not surface as an error until the router kernel failed with an opaque
index-out-of-bounds exception, potentially minutes into a training run. We
replaced this with a hierarchical schema (\texttt{MoEConfig}, six
sub-configs covering model, training, parallelism, checkpointing,
elasticity, and telemetry) with cross-field validators enforcing invariants
such as $\mathrm{top\_k} \le \mathrm{num\_experts}$,
$\mathrm{warmup\_steps} < \mathrm{max\_steps}$, and dtype membership in a
fixed set. Validation failures now report the offending field path and a
human-readable constraint description at load time, before any GPU is
touched.

\subsection{A silent-degradation failure mode, root-caused}
\label{sec:pydantic-bug}

We report the following finding in detail because we consider it
methodologically instructive, not merely a bug-fix footnote. The
configuration module was written with a fallback code path: if the
validation library was not importable, the module would degrade to an
unvalidated pass-through rather than fail to import. The intent was
defensive --- to keep the module importable in minimal environments.

The dependency, however, was never actually added to the package's declared
runtime dependencies. Continuous integration installed the package exactly
as declared, which meant the validation library was never present, which
meant every configuration --- including deliberately invalid ones
constructed by the test suite to verify that validation errors are raised
--- silently passed through the unvalidated fallback. Eight test cases
that asserted a specific exception type would be raised on invalid input
instead observed no exception at all, for an entire release cycle, without
any test failure being reported, because the tests themselves never
exercised the validated code path.

The fix has two parts. First, the validation library is now a declared,
non-optional runtime dependency; second, and more importantly, the
fallback path was deleted entirely. The module now raises an import error
immediately, with installation instructions, if the dependency is absent,
rather than silently substituting a less-safe behavior. We consider this
the correct general policy for any component whose stated purpose is
catching invalid input early: a component that can silently stop
performing its stated function is a more dangerous failure mode than one
that refuses to start.

A second, unrelated defect was found in the same subsystem: the
environment-variable override mechanism used \texttt{yaml.safe\_load} to
parse override values, which is subject to a YAML~1.1 grammar quirk in
which exponential notation without an explicit sign or decimal point
(e.g.\ \texttt{"1e-5"}) parses as a string rather than a float. An override
such as \texttt{MOE\_TRAINING\_\_LEARNING\_RATE=1e-5} would therefore
silently fail with a type error several lines downstream rather than at
the point of parsing. The fix attempts native \texttt{int}/\texttt{float}
coercion before falling back to YAML parsing for booleans, null, and other
scalar types.

\section{Elastic Fault Tolerance and Checkpointing}

The async two-tier checkpointing design (pinned host memory $\rightarrow$
NVMe with \texttt{O\_DIRECT} and atomic rename $\rightarrow$ S3) and the
\texttt{ClusterStateMachine} with round-robin expert resharding are
unchanged from v1/v2 in their core algorithm. One addition was made:
checkpoint metadata now carries an explicit schema version
(\texttt{CHECKPOINT\_SCHEMA\_VERSION}, currently 2), together with the
runtime's own package version and PyTorch version, both read from their
respective single sources of truth rather than hardcoded. On load, a
compatibility check logs at \texttt{info} level for older-but-compatible
checkpoints and at \texttt{warning} level for checkpoints written by a
newer schema than the current runtime understands. This is a diagnostic
aid rather than a hard compatibility gate: missing fields default
gracefully in either direction.

Chaos Scenario B (storage stall) remains fully verified and, as reported
in Section~\ref{sec:eval}, now has real-GPU confirmation. Scenario A (node
kill and recovery) is unchanged at approximately 85\% pass rate, for the
same root cause identified in v2: a Gloo \texttt{connectFullMesh} race
during process-group re-initialization after \texttt{SIGKILL}. Replacing
Gloo with NCCL in the chaos harness remains the top-priority item for v0.4
and requires GPU hardware to implement and validate, which was outside the
scope of the single-GPU validation performed in this cycle.

\section{Observability and MoE-Aware Accounting}

Building on the comm/compute overlap ratio and per-layer expert compute
latency introduced in v2, this release adds four further MoE-specific
telemetry fields, all populated from real per-step measurements rather
than estimated: \texttt{sparse\_mfu} (Model FLOPs Utilization scaled by the
sparse activation fraction $K/E$, which is the theoretically correct
denominator for sparse models and is substantially lower than a naive
dense-equivalent MFU calculation would report); \texttt{dead\_expert\_count}
(the number of experts receiving zero tokens in a given step, an early
warning signal for routing collapse); \texttt{routing\_efficiency} (the
fraction of allocated expert capacity actually consumed); and
\texttt{active\_experts}. A fifth field, \texttt{dropped\_token\_fraction},
reports the fraction of token/slot assignments discarded by the capacity
mechanism of Section~\ref{sec:capacity} and is zero whenever that mechanism
is disabled, which is the default. All five fields are exposed through the
existing four telemetry sinks (structured JSONL, TensorBoard, Prometheus,
and the WandB integration added in v2) without any sink-specific code
duplication, by injecting typed dataclass fields into the existing nested
telemetry dictionaries at construction time.

\section{Implementation and Testing}

The test suite has grown from 148 tests across 33 files at the time of the
v2 preprint to 348 passing tests across 20 files in the fast, CPU-only
tier alone (one additional test is skipped without a CUDA device, and one
is marked as an expected failure for a documented statistical edge case
discussed below), plus a separate multi-process and chaos tier. Three
categories of test infrastructure were added specifically to support this
growth without requiring GPU or multi-process resources for routine
development:

\textbf{Property-based testing.} Nine property-based tests, using the
Hypothesis library, verify invariants --- token conservation, combine-weight
normalization, index validity, complete and non-overlapping expert
ownership under arbitrary $(E, \mathrm{ep\_size})$ combinations, and
configuration cross-field constraints --- across randomly generated inputs
rather than a fixed set of example cases, at fifty generated examples per
property.

\textbf{Mocked collective communication.} A lightweight, thread-based
simulation of \texttt{all\_to\_all\_single} (\texttt{MockDistEnv}) allows
expert-parallel dispatch and combine logic to be exercised at
$\mathrm{ep\_size} > 1$ within a single test process, without
\texttt{torch.distributed} initialization. This is complementary to, not a
replacement for, the existing 2-rank \texttt{mp.spawn} tests, which remain
the source of ground truth for real collective correctness; the mock
exists to make dispatch-logic bugs reproducible in under a second on a
laptop.

\textbf{Dedicated interface test suites.} The high-level router wrapper and
the model registry introduced in this cycle each received a dedicated test
file (33 and 20 tests respectively), separate from the existing kernel-level
and layer-level test suites, following the principle that each public
interface boundary should have tests that would fail if that specific
boundary's contract were violated, independent of whether the underlying
implementation also has coverage.

One previously known issue is worth updating: a stochastic routing-quality
test (\texttt{test\_uniform\_init\_lower\_imbalance}, parametrized over five
random seeds) was reported in v2's antecedent work as intermittently
failing at seed 2. We have since determined this is a genuine statistical
edge case --- at that specific seed, a near-zero-variance gate
initialization produces marginally higher load imbalance (1.625) than a
sharp initialization (1.594) at $E{=}16$, due to random token co-location,
rather than a defect in the load-imbalance calculation itself. The test is
now explicitly marked as an expected failure with the mechanism documented
inline, rather than either silently retrying or being deleted.

\section{Preliminary Evaluation}
\label{sec:eval}

\subsection{Real GPU validation: NVIDIA T4}

This section reports the results the v2 preprint explicitly declined to
report, obtained on a single NVIDIA T4 (16~GB HBM2, 65 TFLOPS FP32) via
Google Colab. We reproduce the exact figures rather than round or restate
them, since the reproduction commands accompanying each are part of the
project's evaluation methodology.

Token conservation held with zero violations across 100 random seeds on
CUDA, matching the CPU reference path's zero-violation result
(N=512, H=128, E=32, K=2). Chaos Scenario B achieved a 10-out-of-10 (100\%)
pass rate under the same storage-stall injection previously validated only
on CPU.

Router kernel throughput, measured with the Triton forward-only kernel
against the fp64 CPU reference path at four configurations, is shown in
Table~\ref{tab:speedup}. The GPU/CPU speedup ratio increases
super-linearly with problem size, from 2.9$\times$ at $N{=}512$ to
80.1$\times$ at $N{=}4096, H{=}2048$ --- the configuration most
representative of production MoE hidden dimensions and expert counts. We
attribute the super-linear scaling to the single-HBM-pass Triton kernel's
bandwidth utilization improving with working-set size, while the CPU fp64
reference path, which is single-threaded and unvectorized, has no
corresponding scaling advantage.

\begin{table}[h]
\centering
\small
\caption{Router forward throughput: T4 GPU vs.\ CPU reference}
\label{tab:speedup}
\begin{tabular}{@{}rrrrrrr@{}}
\toprule
$N$ & $H$ & $E$ & $K$ & CPU (M/s) & GPU (M/s) & Speedup \\
\midrule
512  & 256  & 16 & 2 & 0.747 & 2.141 & 2.9$\times$ \\
1024 & 512  & 32 & 2 & 0.421 & 3.644 & 8.7$\times$ \\
2048 & 1024 & 64 & 2 & 0.236 & 4.832 & 20.4$\times$ \\
4096 & 2048 & 64 & 4 & 0.056 & 4.454 & \textbf{80.1$\times$} \\
\bottomrule
\end{tabular}
\end{table}

We emphasize what these measurements do and do not establish. They confirm
kernel correctness and relative CPU/GPU scaling behavior on real hardware
for the first time. They do not constitute a Model FLOPs Utilization
result at production scale: the T4 is an inference-oriented card with a
BF16 peak roughly $15\times$ lower than an H100 SXM5, and the measured MFU
at the smoke-test batch sizes used is on the order of $10^{-6}$, which is
expected and uninformative at this scale --- it reflects batch size, not
hardware capability. We report it for completeness and transparency rather
than because it is a meaningful result in itself.

\subsection{MoE layer overhead}

Full \texttt{DistributedMoELayer} forward latency was measured against a
single dense SwiGLU FFN baseline (equivalent to $E{=}1, K{=}1$, no router,
no dispatch) at matched $(B, S, H, F)$. At the small batch sizes tractable
for exhaustive small-scale measurement (B$\le$4, S$\le$32), routing and
dispatch overhead on the T4 ranges from 21.8$\times$ to 26.1$\times$
relative to the dense baseline; on CPU, the corresponding range is
9.1$\times$ to 16.2$\times$. We note explicitly that this overhead figure
is dominated by kernel-launch latency at these batch sizes and by the fact
that the dense baseline holds an order of magnitude fewer parameters than
the MoE configuration it is compared against --- the two models are not
capacity-equivalent, and the comparison should be read as a latency-floor
measurement, not a throughput-per-parameter comparison.

\subsection{Comparison to v2's evaluation posture}

Version 2 stated: ``we deliberately do not present large-scale GPU scaling
curves or MFU numbers in this preprint because sustained, reproducible
cluster data is not yet available.'' That statement remains true for
multi-node data. It is no longer true for single-GPU data, which this
section reports. We consider the distinction important: closing the
single-GPU gap was tractable within the resources available for this
release; closing the multi-node gap requires sustained access to a GPU
cluster, which remains the primary blocker for v0.4 and is discussed in
Section~\ref{sec:limitations}.

\section{Discussion, Limitations, and Roadmap}
\label{sec:limitations}

\subsection{Current Maturity (v0.3.3)}

Correctness infrastructure is materially more mature than at v2: the
configuration system now fails loudly rather than silently on invalid
input (Section~\ref{sec:pydantic-bug}); the distributed package is
decomposed into reviewable units; the test suite exercises probabilistic
invariants via property-based testing in addition to fixed examples; and
the router kernel has been validated on real GPU hardware for the first
time, surfacing and fixing a defect that had been latent since v1.

\subsection{Remaining Gaps}

Honestly stated, in the same spirit as v1 and v2:

\begin{itemize}[leftmargin=1.5em]
\item \textbf{Multi-node GPU data.} No benchmark or MFU figures exist
  beyond a single T4. This is the direct continuation of v2's stated P0
  gap and remains the highest-priority item for v0.4.
\item \textbf{Chaos Scenario A.} Unchanged at approximately 85\% pass
  rate; the Gloo-to-NCCL migration identified as the fix in v2 has not
  been implemented, because it requires GPU hardware to validate and was
  outside this cycle's scope.
\item \textbf{Expert capacity mechanism at scale.} The dropping mechanism
  of Section~\ref{sec:capacity} is validated for correctness at toy
  dimensions; its behavior at $E{=}256$ on real hardware --- memory
  footprint, dispatch bandwidth under enforced capacity, and interaction
  with the aggregate load-balancing loss --- has not been measured.
\item \textbf{Capacity re-routing.} The current mechanism drops overflow
  tokens; routing them to a next-best available expert instead (as
  opposed to discarding them) is deferred, because tuning a sensible
  fallback policy requires real expert-parallel bandwidth data that is
  not yet available.
\item \textbf{Pipeline Parallelism under chaos.} The multi-process 1F1B
  schedule has not been exercised under node-failure conditions; the
  existing chaos harness uses a non-pipelined model.
\item \textbf{Non-divisible sequence lengths.} Sequence Parallelism
  still requires $S \bmod \mathrm{tp\_size} = 0$.
\end{itemize}

\subsection{Roadmap}

v0.4 priorities, in order: (1) sustained multi-GPU benchmark data and MFU
validation, requiring cluster access; (2) the NCCL chaos-harness migration;
(3) integrating Pipeline Parallelism into chaos scenarios; (4) hardware
validation of fine-grained expert configurations; (5) Nsight/CUPTI
integration for roofline analysis, which also requires GPU access beyond
what was available for this cycle.

\section{Conclusion}

This preprint reports the closing of a gap explicitly left open by its
predecessor: moe-engine's router kernel and chaos-recovery mechanism have
now been validated on real GPU hardware, at single-device scale, with
reproducible measurements rather than illustrative estimates. In the
process, this validation surfaced a genuine defect --- a missing
\texttt{constexpr} annotation invisible to any CPU-only test --- which
underscores a general point about systems that maintain hardware-specific
code behind a CPU fallback: the fallback's existence is precisely what
allows a hardware-only defect to survive indefinitely in continuous
integration.

In parallel, we undertook a structural refactoring motivated by external
technical review, and in doing so discovered and root-caused a second,
unrelated defect of independent interest: a configuration validator that
had been silently non-functional for an entire release cycle due to an
undeclared dependency, a failure mode we believe generalizes to any system
that treats validated-but-optional as an acceptable combination for
safety-critical logic. We have removed that combination from this codebase
entirely rather than patching around it.

The project's multi-node evaluation gap, identified as the top priority in
v2, remains open in this release for the same reason it was open then:
sustained cluster access. We continue to believe that reporting exactly
what has and has not been measured, at whatever scale is actually
available, is more useful to the field than extrapolating from small-scale
results --- a position this preprint's own single-GPU results section is
itself an instance of.

\section*{Acknowledgments}

The author thanks the PyTorch, Triton, and distributed systems communities
for the foundational tools, the maintainers of Pydantic and Hypothesis for
tools that materially improved this release's correctness posture, and the
many researchers who have openly shared post-mortems and lessons from
large-scale training failures.

\begin{thebibliography}{9}

\bibitem{shoeybi2019megatron}
M.~Shoeybi, M.~Patwary, R.~Puri, P.~LeGresley, J.~Casper, and B.~Catanzaro,
``Megatron-LM: Training multi-billion parameter language models using
model parallelism,'' \emph{arXiv preprint arXiv:1909.08053}, 2019.

\bibitem{rasley2020deepspeed}
J.~Rasley, S.~Rajbhandari, O.~Ruwase, and Y.~He, ``DeepSpeed: System
optimizations enable training deep learning models with over 100 billion
parameters,'' in \emph{Proc. 26th ACM SIGKDD}, 2020.

\bibitem{fedus2021switch}
W.~Fedus, B.~Zoph, and N.~Shazeer, ``Switch transformers: Scaling to
trillion parameter models with simple and efficient sparsity,'' \emph{arXiv
preprint arXiv:2101.03961}, 2021.

\bibitem{hwang2023tutel}
C.~Hwang, W.~Cui, Y.~Xiong, Z.~Yang, Z.~Liu, H.~Hu, Z.~Wang, R.~Salas,
J.~Jose, P.~Ram, J.~Chau, P.~Cheng, F.~Yang, M.~Yang, and Y.~Xiong, ``Tutel:
Adaptive mixture-of-experts at scale,'' \emph{arXiv preprint
arXiv:2206.03382}, 2023.

\bibitem{xmoe2024}
``X-MoE: Scaling Mixture-of-Experts training,'' technical report, 2024.

\bibitem{megascale2025}
``MegaScale-MoE: Large-scale Mixture-of-Experts training infrastructure,''
technical report, 2025.

\bibitem{wu2023lazarus}
Y.~Wu et al., ``Lazarus: Resilient and elastic training of Mixture-of-Experts
models,'' \emph{arXiv preprint}, 2023.

\bibitem{hu2023reliability}
S.~Hu et al., ``Revisiting reliability in large-scale machine learning
research clusters,'' \emph{arXiv preprint}, 2023.

\bibitem{lepikhin2020gshard}
D.~Lepikhin, H.~Lee, Y.~Xu, D.~Chen, O.~Firat, Y.~Huang, M.~Krikun,
N.~Shazeer, and Z.~Chen, ``GShard: Scaling giant models with conditional
computation and automatic sharding,'' \emph{arXiv preprint
arXiv:2006.16668}, 2020.

\end{thebibliography}

\vspace{1em}
\noindent\rule{\textwidth}{0.4pt}
\small
\noindent moe-engine Preprint v3 --- July 2026. The latest version of the
paper and source code are available at
\url{https://github.com/Mattral/Composed-Mixture-of-Experts-Engine}. v1 is
archived at \url{https://doi.org/10.5281/zenodo.20647577}; v2 is archived at
\url{https://doi.org/10.5281/zenodo.20688837}.

\end{document}
